\chapterstart{Структурированная проверка начальных условий}
\label{sec:robust}
\subsection{Что изменяется, а что остаётся фиксированным}
После анализа обучения проведена завершённая диагностическая проверка на этапе разработки (\emph{structured robustness}). В ней нет новых обучений. Оцениваются шесть закреплённых лучших контрольных точек исходной основной серии: SAC-on и TQC-on с seed 0–2, а также прежний классический контроллер. Поколения SAC выбраны на 70/50/80 тыс. переходах, TQC — на 80/40/70 тыс.. Контрольные суммы весов и номинальной траектории закреплены до новых исходов (Приложение~\ref{app:repro}).

В этой главе \emph{off/on означает только режим исполнения одной политики}, обученной с фильтром. SAC-off из предыдущей главы в матрицу не входит. DQN и DDPG также не включены: это ограниченный протокол диагностики выбранных исходных SAC/TQC и классики, а не полное повторное сравнение всех пяти методов. Их robustness-результаты из этого набора вывести нельзя.

Семь групп содержат по 20 общих состояний, всего 140 уникальных стартов. Для каждого контроллера и режима используются те же точные координаты. Матрица имеет $7$ контроллеров $\times2$ режима $\times7$ групп $=98$ заданий по 20 эпизодов, то есть 1960 эпизодов и 980 полных пар. Отдельные 40 timing-эпизодов служили техническому замеру и исключены из научных знаменателей. Название \texttt{main} внутри экспорта robustness обозначает эту оценочную матрицу, а не основную серию обучения.

\subsection{Дизайн стартов и отношение к распределению обучения}
В Таблице~\ref{tab:groups} заданы все группы. Используются $B=\sqrt{2AL}\approx1{,}385638$ м/с — масштаб центральной скорости торможения — и $N=\sqrt{g/l}\approx7{,}378648$ рад/с — естественный частотный масштаб. Это мотивированные моделью масштабы дизайна, а не экспериментально доказанная область успешного подъёма.
\begin{table}[htbp]\caption{Точные правила семи групп, по 20 состояний}\label{tab:groups}
\begin{smalltable}\begin{tabularx}{\linewidth}{@{}p{27mm}X p{40mm}@{}}\toprule
Группа & Значения и генерация & Отношение к области обучающих стартов\\\midrule
Nominal & Независимые $x,v\in[-0{,}02,0{,}02]$, $\theta,\omega\in[-0{,}05,0{,}05]$, равномерно & То же распределение, новые 20 точек\\
Position & $x=\pm\{0{,}02;0{,}04;0{,}06;0{,}08;0{,}10;0{,}12;0{,}15;0{,}18;0{,}21;0{,}23\}$ м & Смешанная: 2 точки на границе носителя, 18 вне по одной оси\\
Velocity & $v=\pm0{,}95B\,k/10$, $k=1,\ldots,10$; максимум $1{,}316356$ м/с & Одноосевое внераспределительное отклонение\\
Angle & $\theta=\pm k\pi/11$, $k=1,\ldots,10$; максимум $2{,}855993$ рад & Одноосевое отклонение; верх точно не включён\\
Angular velocity & $\omega=\pm Nk/10$, $k=1,\ldots,10$ & Одноосевое отклонение\\
Boundary & $x=\pm a$, $a\in\{0{,}20;0{,}23\}$ м; $v=\operatorname{sign}(x)q\sqrt{2A(L-a)}$, $q\in\{0{,}25;0{,}5;0{,}75;0{,}9;0{,}99\}$ & Движение наружу у границы; вне прямоугольника по $x,v$\\
Joint & $s_k=(k/10)\operatorname{diag}(0{,}18;0{,}6;\pi;N)z_k$, $z_k\sim U([-1,1]^4)$, и $-s_k$ & Совместное отклонение по нескольким координатам\\\bottomrule
\end{tabularx}\end{smalltable}
\source{Зафиксированный дизайн экспериментов автора. В одноосевых группах остальные координаты нулевые; в Boundary $\theta=\omega=0$. Точный список — в электронном приложении.}
\end{table}

Случайные значения создавались единственным PCG64 seed 2026091601: сначала 20 nominal-векторов, затем 10 joint-векторов; сетки случайный поток не расходуют. Joint включает зеркальные пары и растущий масштаб, поэтому это не независимая равномерная выборка из большого прямоугольника. Все координаты в SI, угол отсчитывается от нижней вертикали и в исходных данных не оборачивается.

Все 140 стартов удовлетворяют~\eqref{eq:K}. У Boundary $q<1$ оставляет тормозной запас. Для Joint даже верхняя оценка $0{,}18+0{,}6^2/(2A)=0{,}225<L$ гарантирует выполнение позиционных условий отбора. Отбрасывания по результатам контроллера и подрезания координат нет. Тем не менее выбор только $K$ ограничивает область сравнения: не исследуются старты, допустимые для off-reset по рабочему прямоугольнику, но отвергаемые on-reset.

Носитель (support) распределения~\eqref{eq:reset} — совместный четырёхмерный прямоугольник, а не четыре независимые справки о диапазонах. Две Position-точки $x=\pm0{,}02$, остальные координаты нуль, лежат на границе его замыкания, но детерминированная осевая пара не имеет того же распределения, что обучающие старты. Остальные 18 Position и все Velocity, Angle, Angular velocity выходят за этот прямоугольник по одной оси. Все Boundary выходят по двум осям. Среди Joint две точки выходят по двум координатам, десять — по трём, восемь — по четырём. Это конкретная внераспределительная диагностика, а не проверка глобальной устойчивости равновесия.

\subsection{Сначала исходная задача: новые nominal-состояния}
На новых nominal-стартах в режиме on SAC имеет 60/60 успешных исполнений, TQC — 59/60, классика — 20/20. Каждое RL-семейство здесь использует три политики на \emph{одних и тех же двадцати состояниях}; речь не о 60 независимых физических стартах. У TQC seed 1 один nominal-эпизод неуспешен, поэтому все 140 исполнений семи контроллеров дают 139 успехов.

Это проверка из того же исходного распределения на другом фиксированном наборе. Прежние 60/60 TQC на валидации и новые 59/60 на nominal не противоречат друг другу: validation использована для выбора best, nominal20 — другой набор. Новые nominal-результаты показывают перенос на исследованные близкие старты, но не гарантируют успех на всех точках даже малого прямоугольника.
\vref{01_nominal_controllers}{Видеоприложение 1} показывает три выбранных контроллера на общем nominal\_00; это иллюстрация успешного движения, а не все двадцать стартов.

\subsection{Где проявляются ограничения захвата}
Рисунок~\ref{fig:robust_matrix} сохраняет все семь контроллеров и семь групп. Он отвечает на вопрос, одинаково ли меняется качество при разных видах отклонения. Число успешных on-эпизодов из 140 в каждой группе: Nominal 139, Position 122, Velocity 109, Angle 102, Angular velocity 111, Boundary 97, Joint 116. Наиболее трудной по этому равновзвешенному дизайну оказывается Boundary; это свойство выбранных точек и политик, а не универсальный порядок трудности всех возможных возмущений.
\fig{robustness_matrix}{Полная матрица успеха фиксированных контроллеров: каждое число из 20 общих состояний группы. Off/on — исполнение одних весов. Таблица цветов дополнена числами для чтения в чёрно-белой печати.}{1}{fig:robust_matrix}

На всех 140 стартах SAC-on seeds 0/1/2 имеет 135/129/128 успехов, TQC-on — 127/93/99, классика — 85. Таким образом, одинаковый полный успех на старом validation не устранил различия между значениями seed обучения при расширении стартов. По семействам успех off $\to$ on меняется так: SAC — 279/420 $\to$ 392/420, TQC — 101/420 $\to$ 319/420, классика — 61/140 $\to$ 85/140. Знаменатели соответствуют числу исполнений, а не числу независимых обучений. Не следует скрывать неоднородность TQC его средним или трактовать эти доли как вероятности реального применения.

Включение защиты помогает не монотонно на каждом старте. Например, у SAC seed 2 в группе Angular velocity успех уменьшается с 18/20 до 16/20. У классики Boundary остаётся 10/20 в обоих режимах. Наличие фильтра само по себе не обеспечивает новый механизм захвата: он меняет вход тележки, вслед за которым меняется и движение маятника. В этом наборе оценивается прежний классический план с однократным переключением, без перепланирования под каждый старт.
\vref{05_robustness_groups}{Видеоприложение 5} показывает по одному выбранному старту семи групп для SAC seed 0, а
\vref{06_final_controllers}{видеоприложение 6} — все семь контроллеров на joint\_00. Все эти выбранные примеры успешны; неоднородность полной серии видна из таблиц и Рисунка~\ref{fig:robust_matrix}.

\subsection{Парное вмешательство: безопасность и задача}
В Таблице~\ref{tab:robust_total} четыре разных результата оставлены раздельными. Суммарный успех увеличивается с 441/980 до 796/980, то есть с 45,00\% до 81,22\% (на 36,22 процентного пункта). Успех без превышения с допуском возрастает с 394/980 до 796/980, на 41,02 п.п. Разница между 441 и 394 — 47 успешных off-эпизодов, в которых всё же превышалась желаемая граница.
\begin{table}[htbp]\caption{Robustness: исходы по всей фиксированной матрице}\label{tab:robust_total}
\begin{smalltable}\begin{tabularx}{\linewidth}{@{}Xrr@{}}\toprule
Показатель & Off, из 980 & On, из 980\\\midrule
Исходный успех задачи & 441 & 796\\
Успех без resolved-превышения & 394 & 796\\
Resolved-превышение $\pm0{,}24$ м & 585 & 0\\
Достижение горизонта & 446 & 980\\
Горизонт без заключительного удержания & 5 & 184\\
Событие \texttt{position\_limit} на $\pm0{,}25$ м & 534 & 0\\
Событие \texttt{velocity\_limit} & 0 & 0\\\bottomrule
\end{tabularx}\end{smalltable}
\source{Результаты вычислительных экспериментов автора. Числа строгих превышений и превышений с допуском здесь совпадают; определения различаются.}
\end{table}

Во всех 980 on-эпизодах достигнут горизонт, превышений желаемой границы с допуском нет. Однако 184/980 эпизода (18,78\%) не выполняют заключительное удержание. Сохранение тележки в области само по себе не решает задачу подъёма и захвата. Из этих агрегатов нельзя приписать всем 184 ту же причину, что установлена для отдельных DDPG-траекторий в предыдущей главе: DDPG вообще не входит в эту robustness-матрицу.

В 980 парных исходах оба режима успешны 434 раза, только on — 362, только off — 7, ни один — 177. Во всех семи off-only успехах есть resolved-превышение желаемой границы. Поэтому они не являются примерами утраты \emph{успеха без превышения}, хотя являются реальной утратой исходного успеха. Рисунок~\ref{fig:robust_outcomes} показывает обе стороны этого результата.
\vref{03_filter_safety_rescue}{Видеоприложение 3} иллюстрирует устранение остановки по положению;
\vref{04_success_vs_safety}{видеоприложение 4} — успех off с превышением желаемой границы и отсутствие успеха on без превышения. Оба примера выбраны по зафиксированному правилу, а не являются случайной выборкой.
\fig{robustness_outcomes}{Слева — отдельные метрики по 980 эпизодам каждого режима; справа — взаимоисключающие исходы 980 пар. Части слева пересекаются и не должны суммироваться до 100\%.}{1}{fig:robust_outcomes}

Сравнение условно причинно интерпретируемо как эффект включения \emph{всего слоя исполнения} на данных фиксированных весах и стартах: checkpoints, физика и критерии не заменялись. На уровне программного интерфейса меняются также допуск начального состояния и возможность терминального отказа фильтра. Фактически все 140 стартов были допущены в обоих режимах, а отказов фильтра в 1960 эпизодах не было. Поэтому различие состава начальных состояний и новый тип отказа не создают наблюдённый разрыв исходов в этих парах. Реальные прекращения, однако, различаются: 534 off-эпизода остановились по позиции, все on дошли до горизонта. Это часть эффекта исполнения и не основание сравнивать возврат как будто все траектории имели одинаковую длительность.

Общие доли используют равные веса семи групп и веса контроллеров $3$ SAC $+$ $3$ TQC $+$ $1$ классика. Это веса исследовательского дизайна. Сетки и зеркальные пары не образуют независимую случайную выборку из всего пространства, поэтому p-values, доверительные интервалы глобального успеха и перенос на практические частоты стартов не приводятся.

\subsection{Насколько часто фильтр меняет запрос}
Таблица~\ref{tab:interventions} агрегирует вмешательства по числу принятых решений. Всего on изменил 115\,355 из 980\,000 решений, то есть 11,7709\%; отказов нет. При off показатель неприменим. Небольшая доля не означает слабого эффекта: одно раннее торможение может изменить все последующие состояния и обратную связь. Но по одному агрегату нельзя установить момент, ответственный за конкретный исход.
\begin{table}[htbp]\caption{Коррекции запросов в robustness, режим on}\label{tab:interventions}
\begin{smalltable}\begin{tabularx}{\linewidth}{@{}Xrrr@{}}\toprule
Семейство & Изменено / принято & Доля, \% & Средняя / max $|\Delta u|$, м/с$^2$\\\midrule
SAC, три seed & 4209 / 420000 & 1,00 & 2,30 / 7,99\\
TQC, три seed & 57179 / 420000 & 13,61 & 4,65 / 8,00\\
Классика & 53967 / 140000 & 38,55 & 74,70 / 394,64\\\bottomrule
\end{tabularx}\end{smalltable}
\source{Результаты вычислительных экспериментов автора. Средняя величина коррекции рассчитана только среди вмешательств, с суммированием числителей.}
\end{table}

В отличие от запроса обучаемой политики, исходный запрос классики может превышать предел ускорения; величина коррекции включает приведение амплитуды к $\pm4$ м/с$^2$. Поэтому значения 74,70 и 394,64 не означают приложение таких ускорений и не измеряют исключительно дополнительное торможение по $x/v$. Сравнивать их с SAC/TQC как прямую меру «худшего качества защиты» неправильно.

Для классики сохраняется известное расхождение описания и исполнения (\path{known_mismatch}). Спецификация приписывает ей контракт среды обучения SAC, а фактические журналы указывают \path{CourseworkCartPole-v0} без \path{training_contract}. Принятый аудит полных журналов отдельно сверил физику, интегратор, R1 и удержание. Поэтому эти физические метрики сопоставимы, но тождественность программного интерфейса, наблюдений и запросов не установлена. Расхождение ограничивает выводы о полной идентичности процедур, не отменяя измеренных физических результатов.

Завершённая серия показывает перенос на новые nominal-старты и конкретные пределы захвата при структурированных отклонениях. При robustness-проверке веса не изменялись: те же политики оценивались при других начальных условиях. Наблюдаемое снижение успеха относительно nominal характеризует их поведение на выбранных группах; этот эксперимент не устанавливает причины ограниченного обобщения. Это уже полученная development-диагностика; если её используют для изменения метода, она не станет независимым финальным тестом нового метода.
